% Section 5.4 — numbers from paper/stats_study2.py (paper/quant_study2/stats.txt):
%   Study 2 model     eval14p-mpcc-pooled8all-c70a5dd-s42-20260911-1709-ca2b4a7
%   Study 2 baseline  eval14p-baseline-pooled8all-c70a5dd-s42-20260911-1709-ca2b4a7
%   Study 1 reference mpcc-full-perpid-s42-20260909-1257-22c349c (Figure 3, Tables 1-2)
% Figure 5: paper/figures/steering_law_budget_study2.pdf

\subsection{Generalization to Unseen Users}
\label{sec:study2}

Sections~\ref{sec:time} and~\ref{sec:traj} fit one model per person. A
designer evaluating an interface has no such data, so we ask whether a
single model fitted on the Study~1 participants predicts the behavior of
the 14 Study~2 participants, whom it has never seen. A pooled model carries
one lookahead budget, one cutting fraction, and one pace law, so it can
represent a typical user but not the spread between users. We therefore
compare it against population-level quantities first, condition means and
the Steering Law relation, and then against each participant's own trials.

\paragraph{Timing.}
Figure~\ref{fig:study2_law} plots mean completion time per Study~2
condition against $\mathrm{ID}_\kappa$, with $\lambda$ refit on the Study~2
human condition means as in Section~\ref{sec:time}. Study~2 humans follow
$\mathrm{MT} = 0.02 + 0.070\,\mathrm{ID}$ ($R^2 = 0.98$) with
$\lambda = 0.17$~m/rad, below the 0.22~m/rad of Study~1: a unit of turning
costs these participants less, relative to width, than the Study~1
participants. On Study~1's index ($\lambda = 0.22$) the Study~2 human slope
is 0.058 against 0.059 in Study~1, so the two samples agree in the rate at
which time grows with difficulty; the Study~2 intercept is 0.21~s lower.
The pooled model gives $\mathrm{MT} = 0.56 + 0.064\,\mathrm{ID}$
($R^2 = 0.96$), a slope within 9\% of the Study~2 human slope, and a mean
absolute error of 0.41~s (10\% of the 4.1~s mean human time) across the 25
conditions. The pooled baseline gives $\mathrm{MT} = 5.75 + 0.052\,\mathrm{ID}$
($R^2 = 0.12$) with a mean absolute error of 5.8~s: its completion time
barely depends on difficulty. The model's intercept lies 0.54~s above the
human one, the same offset as in Study~1 (0.60~s), while its slope is
closer to the human slope than in Study~1 (9\% against 16\%). The constant
offset is therefore a property of the model and not of the new sample. It
is largest in gentle sinusoidal and straight corridors (model over human
time 1.17 and 1.15) and absent in sharp sinusoidal corridors (0.99): the
model is slow where the corridor leaves pace free.
By width, the model is closest at the three widths absent from Study~1,
20, 30, and 40~mm (1.04 to 1.09), and farthest at the two it was fitted
on, 10 and 50~mm (1.12 and 1.10): the pace law interpolates.

\paragraph{Spread.}
Across the 14 participants, the SD of completion time per condition is
1.53~s (from 0.96~s in straight to 2.17~s in corner corridors). The pooled
model's run-to-run SD, which reflects only execution noise and latency
sampling, is 0.19~s. The model under-spreads by a factor of 8, as expected:
the between-participant variance that the per-participant models of
Study~1 reproduced (1.19~s against 1.63~s across the 8 participants) came
from the fitted lookahead budget, cutting fraction, and pace parameters,
which the pooled model collapses to one value each. The model's run-to-run
SD is also below the participants' own round-to-round SD of 0.61~s.

\paragraph{Trajectories.}
Table~\ref{tab:study2_metrics} reports the four metrics of
Section~\ref{sec:traj} for the pooled model and the pooled baseline on
Study~2, computed per participant--condition pair as in
Table~\ref{tab:overall_metrics}. The pooled model has a lateral RMSE of
3.2~mm, a speed RMSE of 0.109~m/s, a speed correlation of 0.22, and a
time difference of 36.5\%. As a reference, the last row of the table
repeats the per-participant models' overall values from Study~1. Speed
correlation, which measures where along the path the model slows and
speeds up, is at the same level (0.22 against 0.21). Lateral RMSE and
time difference are higher than the reference values (2.5~mm and 17.2\%),
as they are measured against each individual's trials while the pooled
model predicts one typical trajectory: the 10\% timing error of the
condition means above becomes 36\% at the level of individual
participants, and the difference is the individual spread. Lateral RMSE
grows with width, from 1.1~mm at 10~mm to 5.2~mm at 50~mm, since wide
corridors leave room for individual paths that a single cutting fraction
cannot follow; corners carry the largest lateral error (4.1~mm). Timing
error is largest in straight (53\%) and gentle sinusoidal (42\%)
corridors, where the corridor leaves pace to the person, and smallest in
sharp sinusoidal corridors (24\%), where the corridor dictates it. The
pooled baseline is worse on every metric (4.2~mm, 0.213~m/s, 0.06,
203\%). Paired one-sided Wilcoxon signed-rank tests over participants
($n = 14$) confirm the difference: the model has lower lateral RMSE, speed
RMSE, and time difference for every participant ($W = 0$, $p < 0.001$
each) and higher speed correlation for 11 of 14 ($p = 0.007$). No
simulated run of either model timed out or diverged in Study~2.

What transfers is the dependence of pace on width and curvature: $\gamma$,
the pace law, and the budget mechanism are population-level and carry the
Steering Law slope to new users, including corridor widths the model was
never fitted on. What does not transfer is the personal cutting fraction
and lookahead budget, which set corner shape and the spread of times.
Individual fitting therefore remains necessary to predict a specific user,
while the pooled model suffices to predict how long a typical user takes
and where along the path they slow down.

\begin{figure}[t]
    \centering
    \includegraphics[width=0.6\linewidth]{figures/steering_law_budget_study2.pdf}
    \caption{Movement time against the curvature-aware index of difficulty
    for the 14 Study~2 participants (blue), the pooled model (orange), and
    the pooled baseline (gray), one point per task condition averaged across
    participants, with linear fits; $\lambda = 0.17$~m/rad fit on the
    Study~2 human means. The pooled model recovers the Steering Law of
    users it has never seen with a slope within 9\% of theirs.}
    \Description{Scatter plot of movement time in seconds, from 0 to about
    27, against curvature-aware index of difficulty from 0 to 190. Blue
    circles for humans and orange triangles for the model lie close to each
    other along a rising line from about 1 s at ID 10 to about 13 s at ID
    190; their fitted lines nearly coincide, the orange line slightly above
    the blue at low ID and slightly below at high ID. Gray squares for the
    baseline scatter widely between 1 and 18 s with no clear trend and a
    shallow gray fit line starting near 6 s. The legend gives the fitted
    equations: Human MT = 0.02 + 0.070 ID with R squared 0.98, Ours MT =
    0.56 + 0.064 ID with R squared 0.96, Baseline MT = 5.75 + 0.052 ID with
    R squared 0.12, and lambda = 0.17 m/rad fit on human.}
    \label{fig:study2_law}
\end{figure}

\begin{table}[t]
    \caption{Trajectory metrics on Study~2 by corridor type (mean $\pm$ SD
    across participant--condition pairs, 14 participants $\times$ 25
    conditions): the single pooled model and the pooled baseline, both
    fitted on the 8 Study~1 participants. The last row repeats, as a
    reference, the per-participant models' overall values on Study~1 from
    Table~\ref{tab:overall_metrics}. Lateral and speed RMSE are relative to
    the participant's own trajectory after resampling by progress; time
    difference is $|t_{\mathrm{model}} - t_{\mathrm{human}}| / t_{\mathrm{human}}$.}
    \Description{Table with six corridor-type blocks, straight, corner,
    sinusoidal, gentle sinusoidal, sharp sinusoidal, and all, each with a
    pooled-model row and a pooled-baseline row, and a final reference row
    for the per-participant models of Study 1. Four columns give lateral
    RMSE in millimetres, speed RMSE in metres per second, speed
    correlation, and time difference in percent. In the all block the
    pooled model scores 3.2 mm, 0.109 m/s, 0.22, and 36.5 percent, the
    pooled baseline 4.2 mm, 0.213 m/s, 0.06, and 203 percent, and the Study
    1 reference 2.5 mm, 0.087 m/s, 0.21, and 17.2 percent. The baseline is
    worse than the pooled model on every metric in every block; speed
    correlation of the pooled model matches the Study 1 reference while its
    lateral and time errors are higher.}
    \label{tab:study2_metrics}
    \small
    \setlength{\tabcolsep}{4pt}
    \begin{tabular}{l l cccc}
        \toprule
        Corridor type & Model & Lat.\ (mm) & Spd.\ (m/s) & Spd.\ $r$ & Time (\%) \\
        \midrule
        Straight           & Pooled model               & $2.1 \pm 1.2$ & $0.216 \pm 0.166$ & $0.26 \pm 0.44$ & $53.1 \pm 48.4$ \\
                           & Pooled baseline            & $2.0 \pm 1.1$ & $0.502 \pm 0.187$ & $-0.19 \pm 0.42$ & $38.6 \pm 23.9$ \\
        Corner             & Pooled model               & $4.1 \pm 2.4$ & $0.092 \pm 0.052$ & $0.63 \pm 0.14$ & $32.7 \pm 24.9$ \\
                           & Pooled baseline            & $4.8 \pm 1.3$ & $0.130 \pm 0.069$ & $0.19 \pm 0.22$ & $71.0 \pm 56.7$ \\
        Sinusoidal         & Pooled model               & $3.5 \pm 1.9$ & $0.069 \pm 0.049$ & $0.03 \pm 0.31$ & $30.5 \pm 23.1$ \\
                           & Pooled baseline            & $4.5 \pm 1.4$ & $0.164 \pm 0.104$ & $0.15 \pm 0.31$ & $470.9 \pm 308.2$ \\
        Gentle sinusoidal  & Pooled model               & $3.0 \pm 1.5$ & $0.107 \pm 0.091$ & $-0.04 \pm 0.30$ & $42.3 \pm 34.9$ \\
                           & Pooled baseline            & $4.5 \pm 1.7$ & $0.132 \pm 0.086$ & $0.16 \pm 0.37$ & $78.5 \pm 71.1$ \\
        Sharp sinusoidal   & Pooled model               & $3.5 \pm 1.9$ & $0.060 \pm 0.043$ & $0.22 \pm 0.23$ & $23.9 \pm 14.9$ \\
                           & Pooled baseline            & $5.1 \pm 1.2$ & $0.139 \pm 0.087$ & $0.02 \pm 0.20$ & $356.6 \pm 232.9$ \\
        \midrule
        All                & Pooled model               & $3.2 \pm 1.9$ & $0.109 \pm 0.108$ & $0.22 \pm 0.38$ & $36.5 \pm 33.1$ \\
                           & Pooled baseline            & $4.2 \pm 1.8$ & $0.213 \pm 0.185$ & $0.06 \pm 0.35$ & $203.1 \pm 250.4$ \\
                           & Per-participant, Study 1   & $2.5 \pm 2.4$ & $0.087 \pm 0.110$ & $0.21 \pm 0.31$ & $17.2 \pm 16.5$ \\
        \bottomrule
    \end{tabular}
\end{table}
